\section*{Appendix}
\label{appendix:prompts}

\begin{tcolorbox}[
    colback=lightgray!20, 
    colframe=gray, 
    title={\textbf{Claim Extraction Prompt}}, 
    fonttitle=\bfseries,
    sharp corners=southwest,
    coltitle=black,
    label=prompt:p3
]
\small

You are given a passage of text. Your task is to extract one-sentence standalone factual claims from the passage.

\textbf{Each claim must:}
\begin{itemize}
    \item Be self-contained (i.e., understandable on its own without requiring coreference resolution)
    \item Not rely on ambiguous references like ``this'', ``that'', ``they'', etc.
    \item Be phrased as a complete declarative sentence
    \item Be verifiable (e.g., ``Riot Games has no plans for League of Legends 2.'' is acceptable; ``They have no plans for League of Legends 2'' is not)
\end{itemize}

For each claim, also provide the exact verbatim supporting text span from the original text that led you to extract the claim.

\vspace{0.01cm}
\textbf{Output JSON Keys:}
\begin{itemize}
    \item \texttt{"claim1"}: standalone claim
    \item \texttt{"supporting\_text\_span1"}: verbatim supporting text span from original
    \item \texttt{"claim2"}: standalone claim
    \item \texttt{"supporting\_text\_span2"}: verbatim supporting text span from original
    \item \texttt{...}
\end{itemize}

\vspace{0.01cm}
\textbf{Task:} \newline
Here is the input text: \texttt{\{text\}}

\end{tcolorbox}

\begin{tcolorbox}[
    colback=lightgray!20, 
    colframe=gray, 
    title={\textbf{Temporal-Sequence Multi-Hop QA Prompt}}, 
    fonttitle=\bfseries,
    sharp corners=southwest,
    coltitle=black,
    label=prompt:p4
]
\small

You are generating \texttt{\{n\_pairs\}} \textbf{temporal-sequence} multi-hop QA pairs. \\
Buckets of claims (grouped by doc\_id) are below: \\
\texttt{\{json.dumps(buckets, indent=2)\}}

\vspace{0.2cm}
\textbf{Rules for each QA pair:}
\begin{itemize}
    \item Combine claims from at least \textbf{\{min\_buckets\} different buckets}.
    \item Each chosen claim must mention a year, date, or explicit timeframe.
    \item Ask for an ordering or an interval (e.g., ``How many years after …?'').
    \item Every referenced fact must be essential for the answer.
    \item The question should not contain the answer or steps to get the answer.
    \item The question should not refer to the documents, they should be general.
    \item List every claim you used by its doc\_id and claim\_id.
    \item Provide one concise answer (year, number of years, earlier event, etc.).
\end{itemize}

\vspace{0.2cm}
\textbf{Output format (JSON list):}
\begin{itemize}
    \item \texttt{"used\_claims"}: list of \{\{doc\_id, claim\_id, claim\}\}
    \item \texttt{"question"}: string
    \item \texttt{"answer"}: string
\end{itemize}

Return only the JSON—no extra commentary.

\end{tcolorbox}

\begin{tcolorbox}[
    colback=lightgray!20, 
    colframe=gray, 
    title={\textbf{Comparison-Style Multi-Hop QA Prompt}}, 
    fonttitle=\bfseries,
    sharp corners=southwest,
    coltitle=black,
    label=prompt:p5
]
\small

You are generating \texttt{\{n\_pairs\}} \textbf{comparison-style} multi-hop QA pairs. \\
Buckets of claims (grouped by doc\_id) are below: \\
\texttt{\{json.dumps(buckets, indent=2)\}}

\vspace{0.2cm}
\textbf{Rules for each QA pair:}
\begin{itemize}
    \item Use claims from at least \textbf{\{min\_buckets\} different buckets}.
    \item Frame the question so the solver must contrast values (higher, lower, difference, ratio, earlier, later).
    \item Keep wording tight—only include values essential for the comparison.
    \item The question should not contain the answer or steps to get the answer.
    \item The question should not refer to the documents, they should be general.
    \item List every claim you used by its doc\_id and claim\_id.
    \item Provide a single, concise answer (number, entity, etc.).
\end{itemize}

\vspace{0.2cm}
\textbf{Output format (JSON list):}
\begin{itemize}
    \item \texttt{"used\_claims"}: list of \{\{doc\_id, claim\_id, claim\}\}
    \item \texttt{"question"}: string
    \item \texttt{"answer"}: string
\end{itemize}

Return only the JSON—no extra commentary.

\end{tcolorbox}

\begin{tcolorbox}[
    colback=lightgray!20, 
    colframe=gray, 
    title={\textbf{Conjunction-Style Multi-Hop QA Prompt}}, 
    fonttitle=\bfseries,
    sharp corners=southwest,
    coltitle=black,
    label=prompt:p6
]
\small

You are generating \texttt{\{n\_pairs\}} \textbf{conjunction-style} multi-hop QA pairs. \\
Buckets of claims (grouped by doc\_id) are below: \\
\texttt{\{json.dumps(buckets, indent=2)\}}

\vspace{0.2cm}
\textbf{Rules for each QA pair:}
\begin{itemize}
    \item Combine facts from at least \textbf{\{min\_buckets\} different buckets} (can use more).
    \item The answer must change if any referenced claim were false (logical AND).
    \item No trivia—every fact must be essential.
    \item The question should not contain the answer or steps to get the answer.
    \item The question should not refer to the documents, they should be general.
    \item List every claim you used by its doc\_id and claim\_id.
    \item Provide a single, concise answer (entity, number, date, etc.).
\end{itemize}

\vspace{0.2cm}
\textbf{Output format (JSON list):}
\begin{itemize}
    \item \texttt{"used\_claims"}: list of \{\{doc\_id, claim\_id, claim\}\}
    \item \texttt{"question"}: string
    \item \texttt{"answer"}: string
\end{itemize}

Return only the JSON—no extra commentary.

\end{tcolorbox}
